\section{KẾT LUẬN}

Đề tài đã hoàn thành việc xây dựng và đánh giá thực nghiệm hệ thống tác nhân tự tiến hóa (Self-Evolving AI Agents) hỗ trợ chẩn đoán và xử lý sự cố mạng viễn thông trên nền tảng benchmark NIKA. Thông qua việc tích hợp hai mô-đun cốt lõi là \textbf{Procedural Memory} và \textbf{Tool Refinement}, hệ thống cho thấy năng lực tự cải tiến hành vi vượt trội qua các episode chẩn đoán thực tế. Kết quả thực nghiệm trên các mô hình ngôn ngữ lớn làm backbone gồm GPT-OSS-120B và Qwen3.5-122B-A10B-FP8 chứng minh hướng đi này giúp nâng cao đáng kể độ chính xác xác định nguyên nhân gốc rễ (RCA), đưa tỷ lệ lỗi gọi công cụ về mức tuyệt đối và tối ưu hóa đáng kể số bước suy luận mà không cần tái huấn luyện mô hình. Tính mới của đề tài nằm ở việc định hình kiến trúc tác nhân như một thực thể động có khả năng tự sửa lỗi và tích lũy tri thức có cấu trúc từ các vết lỗi trong quá trình tương tác thực tế với môi trường mạng viễn thông.

Về mặt đóng góp cá nhân, sinh viên đã trực tiếp thực hiện toàn bộ quy trình từ nghiên cứu lý thuyết về các phương pháp tự tiến hóa, thiết kế kiến trúc tích hợp với nền tảng NIKA, hiện thực mã nguồn cho hai mô-đun học kinh nghiệm, thiết lập pipeline thực nghiệm tự động hóa và phân tích số liệu kết quả. Trong các bước phát triển tiếp theo, hệ thống hướng tới việc tối ưu hóa khả năng chọn lọc và pruning bộ nhớ để tránh lãng phí chi phí token, đồng thời nghiên cứu khả năng thích ứng linh hoạt trên nhiều topology mạng thực tế quy mô lớn hơn. Khi được hoàn thiện, giải pháp này hứa hẹn sẽ trở thành một công cụ trợ lý thông minh đắc lực cho các kỹ sư vận hành hệ thống mạng, giúp tự động hóa quy trình thu thập bằng chứng, khoanh vùng sự cố và chuẩn hóa các báo cáo kỹ thuật một cách nhanh chóng và chính xác.
